\lstset{
  language=Python,                	 % elegir el lenguaje del código
  basicstyle=\ttfamily\footnotesize,
  numberstyle=\tiny\color{Output},
  stringstyle=\color{Input},
  keywordstyle=\color{teal},
  commentstyle=\color{codegreen},
  numbers=left,                   % dónde colocar los números de línea
  stepnumber=1,                   % paso entre dos números de línea
  numbersep=5pt,                  % distancia de los números de línea respecto al código
  backgroundcolor=\color{SrcBckgrnd},  % color de fondo. Debe añadirse \usepackage{color}
  showspaces=false,               % mostrar espacios con guiones bajos especiales
  showstringspaces=false,         % subrayar los espacios dentro de las cadenas
  showtabs=false,                 % mostrar tabulaciones dentro de las cadenas con guiones bajos especiales
  tabsize=2,                      % establece el tamaño predeterminado de la tabulación a 2 espacios
  captionpos=b,                   % coloca la leyenda en la parte inferior
  breaklines=true,                % activa el salto automático de líneas
  breakatwhitespace=true,         % los saltos automáticos solo ocurren en espacios en blanco
  title=\lstname,                 % muestra el nombre del archivo incluido con \lstinputlisting
}
